% Hints to the exercises for VI/03.
\begin{enumerate}[label=\textbf{E\arabic*.},leftmargin=*]
\item A polynomial restricts to the zero function exactly when it vanishes at every point.
\item Apply the previous exercise to \(f-g\).
\item Use the fact that a polynomial over an infinite field vanishing on all of \(k^n\) is zero.
\item The equations \(x_i-a_i=0\) force every coordinate.
\item Eliminate \(y\) using \(by=c-ax\).
\item Send \([x]\mapsto t\) and \([y]\mapsto t^2\); construct the inverse by \(t\mapsto[x]\).
\item In the quotient \(xy=1\), so \([x]\) is invertible.  Use the universal description of Laurent polynomials.
\item Every residue class is a polynomial in the coordinate classes; the kernel of the quotient map contains exactly the relations.
\item If \([f]^m=0\), then \(f^m\in I(X)\); use radicality of \(I(X)\).
\item Constants show surjectivity.  Then use \(k[X]/\ker(\operatorname{ev}_a)\cong k\).
\item Start from \(I(X)\subseteq I(Y)\) and use the induced quotient map.
\item Use the Chinese remainder theorem for \((x-a)\) and \((x-b)\).  For idempotents, use Lagrange interpolation.
\item Apply the Chinese remainder theorem to the pairwise comaximal ideals \((x-a_i)\).
\item In \(k[x,y]/(xy)\), inspect the classes of \(x\) and \(y\).
\item The class of \(x\) in \(B\) is nonzero but has square zero.
\item First note \(x^q-x\) vanishes on \(\mathbb F_q\).  Divide any vanishing polynomial by it and use the root bound on the remainder.
\item A function on \(X\times Y\) is just a list of \(rs\) values; polynomial representability identifies the ring with all such lists.
\item Multiply \((\bar x+i\bar y)(\bar x-i\bar y)\) and use \(\bar x^2+\bar y^2=1\).
\end{enumerate}
